\section{AI Assistance and Formal Verification}
\label{sec:ai}
\subsection{A workflow with explicit acceptance conditions}
The research used an ARIS-style workflow: a persistent formal specification, linked literature and experiment records, parallel exploration, and separate review of proposed arguments. The specification fixed the unrestricted matrix quantifiers before search. A positive candidate required full rank and a complete nonzero-message distance check; a negative conclusion required a reduction covering the entire target class. The research memory retained unsuccessful routes with their actual scope, preventing a fixed-seed exclusion from becoming a global nonexistence claim in a later session.

The final proof does not require replaying the discovery trajectory. Candidate generation, constraint-programming models, prompts, and search order guided exploration. The accepted argument uses the reductions of Section~\ref{sec:reduction}, the published class bound, and the finite evidence of Section~\ref{sec:computation}. Figure~\ref{fig:flow} shows this dependency structure.

\begin{figure}[t]
\centering
\begin{tikzpicture}[font=\scriptsize,>=Latex,
 box/.style={draw,rounded corners=2pt,align=center,text width=7.3cm,inner sep=4pt},
 node distance=3.8mm]
\node[box] (a) {Arbitrary $[43,10,\ge17]$ code\\Parity extension; E1; dual weight four; shortening};
\node[box,below=of a] (b) {Full-support $[40,7,18]$ mother\\E2 + 172 actual inequivalent representatives};
\node[box,below=of b] (c) {Two complete scans of $E_{40}/K$\\Every pair in $F_K$ has a checked rejection witness};
\node[box,below=of c] (d) {Main exclusion + three puncturing/shortening corollaries};
\node[box,below=6mm of d] (e) {$[42,16,\ge13]\ \Rightarrow\ [41,14,14]$\\Parity, shortening, normalization; contradiction to E3};
\draw[->] (a)--(b);\draw[->] (b)--(c);\draw[->] (c)--(d);
\end{tikzpicture}
\caption{Two routes to five exclusions. AI assists in finding and checking the routes. Published theorems, mathematical deductions, and explicit computations supply the proof dependencies. The fifth target uses E3 rather than the 172-mother computation.}
\label{fig:flow}
\end{figure}

\subsection{Recorded interventions and their scope}
The project records support two concrete examples of model contribution. A Gemini consultation suggested testing native early-rejection filters. The subsequent profile found that candidates in the tested instance needed about 3.38 distance checks on average before rejection, motivating further local implementation of bit-parallel low-bit masks. The implemented mask algorithm and its full-domain justification were developed in the local research process; the model's performance prediction was revised by measurement. The profile concerns an early fixed subcode, rather than the final 172-case workload.

A Grok review sharpened the missing completeness check. At that stage, independent validation established that saved cosets satisfied the constraints, while the review packet already noted that an independent full scan was missing. The review explicitly required a second implementation to decide acceptance or rejection over the entire quotient and compare complete output sets. The project then implemented the opposite-pivot, Gray-order scan described above. This intervention has a concrete consequence: it tests omitted admissible cosets, which member validation alone cannot test.

Local agents also worked on source retrieval, structural deductions, mother reconstruction, code, and proof exposition. Later consultations checked equations, quantifiers, classification scope, and evidence connections. These are documented research and review contributions. The final verification procedures operate on ordinary mathematical and computational objects and require no model API.

\subsection{What the Lean developments establish}
Both formal projects use Lean 4.34.1 and pinned mathlib. The first formalizes parity extension, puncturing, shortening, the Griesmer argument, the MacWilliams calculation, the full-support mother reduction, quotient normalization, classification coverage, and the corollaries. Actual mother matrices and lower-bound generators are connected to the abstract code definitions and checked by kernel evaluation.

The final main theorem takes four explicit parameters: E1; E2; pairwise coordinate inequivalence of the catalogue; and pair incompatibility for each catalogue member. Its type therefore records the boundary between kernel-checked deduction and the published or external computational facts. A finer interface expands the scan and certificate assumptions. The second project proves the entire reduction to exact $[41,14,14]$ and takes E3 as an explicit parameter.

The audited theorem dependencies contain the usual Lean axioms \texttt{propext}, \texttt{Classical.choice}, and \texttt{Quot.sound}, with no \texttt{sorryAx}. Explicit parameters do not appear in this axiom list. Consequently, a clean axiom audit verifies the deductions but does not discharge E1--E3, catalogue inequivalence, or the exhaustive computations. Appendix~\ref{app:formal} gives the exact interfaces.
